\documentclass{article}
\usepackage[utf8]{inputenc}
\usepackage[T1]{fontenc}
\usepackage{geometry}
\geometry{
    a4paper,
    left=2.5cm,
    right=2.5cm,
    top=2.5cm,
    bottom=2.5cm
}
\usepackage{graphicx}
\usepackage{float}
\usepackage{hyperref}
\usepackage[english]{babel}
\usepackage{amsmath}
\usepackage{amssymb}
\usepackage{array}
\usepackage{booktabs}
\usepackage{caption}
\usepackage{subcaption}
\usepackage{enumitem}
\usepackage{fancyhdr}
\usepackage{titlesec}
\usepackage{url}
\usepackage{textcomp}
\usepackage{xcolor}
\usepackage{listings}
\usepackage{multirow}
\usepackage{longtable}
\usepackage{tabularx}
\usepackage{times}

% Set up page style
\pagestyle{fancy}
\fancyhf{}
\rfoot{\thepage}
\renewcommand{\headrulewidth}{0pt}

% Title formatting
\titleformat{\chapter}[display]{\normalfont\bfseries}{\chaptertitle\thechapter}{0pt}{\Huge}
\titlespacing*{\chapter}{0pt}{50pt}{40pt}

% Document info
\title{Template Document}
\author{Academic Support Agent}
\date{\today}

\begin{document}
\maketitle

%---------------------------------------------------------------
% RESUMEN (edici n: AUTORES. Ir a abstract.tex)
% Content from abstract would be included here

%---------------------------------------------------------------
% CORREO ELECTR NICO DEL AUTOR DE CONTACTO (edici n: AUTORES)
\email{* zgamma@citedef.gob.ar}  
%---------------------------------------------------------------
% FORMATO 1er P GINA (edici n: ANALES AFA)
\thispagestyle{empty}     
%\thispagestyle{firstpage}  

%===============================================================
% CUERPO PRINCIPAL DEL TRABAJO (edici n: AUTORES)
\section{INTRODUCCI N}

Estas instrucciones tienen como objetivo ayudar a los autores que deseen publicar en Anales AFA a preparar sus manuscritos. Una vez que la versi n final del trabajo haya sido aceptada, el autor recibir  la \textit{Prueba de Galera} para su revisi n y aceptaci n final.

Este documento es en s  mismo un ejemplo de c mo quedar n los trabajos en su forma final para publicaci n. Los art culos pueden ser escritos en ESPA OL, PORTUGU S o INGL S. El t tulo, resumen y palabras clave deber n estar escritos tanto en el idioma elegido como en un segundo idioma. Esta pauta tiene el prop sito de cumplir con los requisitos de indexaci n de la revista en distintos cat logos.
Si elije escribir el art culo en espa ol o portugu s, deber  incluir el t tulo, resumen y palabras clave tanto en el idioma elegido y como en INGL S; esto  ltimo es a los fines de su inclusi n en los sistemas internacionales de citaciones. 

Se solicita que todos y cada uno de los autores proporcionen su n mero de identificador normalizado ORCID; para mayor informaci n, ingresar a \urlstyle{same}\url{https://orcid.org/}.


\section{GU A PARA LOS AUTORES}
\subsection{Instrucciones para la escritura}

Los manuscritos deben inscribirse en un rect ngulo de $17 \times \unit[25,2]{cm^{2}}$, como se muestra en la Fig. \ref{fig01}, en papel formato DIN A4 ($21 \times \unit[29,7]{cm^{2}}$) y respetando los m rgenes que se indican. Las ecuaciones, tablas, figuras y referencias deben seguir un esquema num rico secuencial.
% EJEMPLO: FIGURA EN 1 COLUMNA
 \begin{figure}[ht]
    \centering
    \includegraphics[width=0.44\textwidth]{fig1.jpg}
  \caption{Formato de p gina recomendado}
  \label{fig01}
  \end{figure}
Los manuscritos deben enviarse en formato PDF para la revisi n por pares. Nuestro sistema de env o de manuscritos (visitar  \urlstyle{same}\url{https://anales.fisica.org.ar/}) le guiar  paso a paso en el proceso para ingresar los datos de su art culo y cargar los archivos. El autor de correspondencia recibir  informaci n sobre el proceso de revisi n de su manuscrito a trav s del email.  
 
 % EJEMPLO: TABLA QUE ABARCA 2 COLUMNAS
\begin{table*}[ht]
\caption{Estilos recomendados.}
\label{tab_01}
    \centering
    \small
    \begin{tabular}{|m{2cm}|m{1.5cm}|m{7.8cm}|c|}
    \hline
    \centering
    \textbf{Estilo} & \centering\textbf{Tama o de la fuente} & \centering\textbf{Descripci n} & \textbf{Ejemplo} \\ 
    \thickhline
    \centering T tulo General & \centering 14 pt. & Fuente: negrita; sin sangr a; p rrafo centrado; interlineado sencillo; espaciado posterior: 12 pt.  &  \large\textbf{T TULO} \\ 
    \hline
    \centering Autor & \centering 11 pt. &  Fuente: negrita; sin sangr a; p rrafo centrado; interlineado sencillo; espaciado posterior: 12 pt. & \large\textbf{{Autor}} \\ 
    \hline
    \centering Direcci n & \centering 9 pt. & Fuente: it lica; sin sangr a; p rrafo centrado; interlineado sencillo; espaciado anterior y posterior: 0 pt.  &  \small\textit{{Direcci n}}\\
    \hline
    \centering Revisi n y aceptaci n & \centering 9 pt.  & Fuente: normal; p rrafo centrado; interlineado sencillo; espacio entre caracteres expandido a 1 pt.; espaciado anterior: 12 pt; espaciado posterior: 6 pt.  & \small{Aceptado}\\
    \hline
    \centering Resumen & \centering 9 pt. & Fuente: normal; indentado a derecha e izquierda: 1 cm; p rrafo centrado; interlineado sencillo; espaciado anterior: 0 pt.; espaciado posterior: 6 pt.  & \small{Resumen}  \\
    \hline
    \centering Palabras clave  &  \centering 9 pt.  & Fuente: it lica; sin sangr a; justificado a izquierda; interlineado sencillo; espaciado anterior y posterior: 6 pt.  & \small{\textit{Palabras clave}}  \\
    \hline
    \centering Texto  & \centering 10 pt.  & Fuente: normal; sangr a de primera l nea: 0,5 cm; justificado a derecha e izquierda; interlineado sencillo; espaciado anterior y posterior: 0 pt.  & Este es un ejemplo ...  \\
    \hline
    \centering T tulo 1  & \centering 10 pt. &  Fuente: negrita; sin sangr a; justificado a izquierda; interlineado sencillo; espaciado anterior: 6 pt.; espaciado posterior: 3 pt.   &  \textbf{I. SECCI N} \\
    \hline
    \centering T tulo 2 & \centering 10 pt.  & Fuente: negrita; sin sangr a; justificado a izquierda; interlineado sencillo; espaciado anterior: 6 pt.; espaciado posterior: 3 pt.  & \textbf{Subsecci n}  \\
    \hline
    \centering Ecuaciones & \centering 10 pt.  &  Centrado en la columna; puede incluir n mero de ecuaci n justificado a la derecha; espaciado anterior: 0 pt.; espaciado posterior: 3 pt.  &  $p= 2\pi d$ \hspace{0.5cm}  (1)  \\
    \hline
    \centering T tulo de tablas  & \centering 10 pt.  & Fuente: it lica; sin sangr a; p rrafo centrado; interlineado sencillo; espaciado anterior: 6 pt.; espaciado posterior: 9 pt.  &  \textit{TABLA 1}  \\
    \hline
    \centering Descripci n de figuras  & \centering 9 pt.  &  Fuente: it lica; sin sangr a; justificado a izquierda; interlineado sencillo; espaciado anterior: 6 pt.; espaciado posterior: 9 pt.  &  \textit{Figura 1:}  \\
    \hline
    \centering Agradecimientos & \centering 10 pt. &  Fuente: negrita; sin sangr a; justificado a izquierda; interlineado sencillo; espaciado anterior: 6 pt.; espaciado posterior: 3 pt. No numerar  esta secci n y situarla al final del manuscrito (antes de las Referencias). &  \textbf{AGRADECIMIENTOS}\\
    \hline
    \centering Referencias & \centering 9 pt. & Sangr a francesa: 0,5 cm; justificado a derecha e izquierda; interlineado sencillo; espaciado anterior: 0 pt.; espaciado posterior: 3 pt. No numerar esta secci n.   &    \\
    \hline
    \centering Nota al pie  &  \centering 9 pt.  & Fuente: it lica; sin sangr a; justificado a izquierda; interlineado sencillo; espaciado anterior y posterior: 0 pt.  &  \small{\textit{* email}}   \\
    \hline
    \end{tabular}
\end{table*}

\normalsize

La versi n final, ya corregida y aprobada, se entregar  tanto en PDF como en una versi n editable. Para este  ltimo se aceptar n archivos en Word (extensi n .doc o .docx) o \LaTeX{ }  (archivos del manuscrito .tex y, en caso que corresponda, archivo .bib). Adem s, se solicita a los autores adjuntar las correspondientes figuras en formato JPG   PNG en ambos casos. Los archivos deber n ser enviados por correo electr nico a la direcci n \textcolor{blue}{analesafaba@gmail.com}  y con copia a la direcci n \textcolor{blue}{analesafatan@gmail.com}.

El presente documento, as  como las instrucciones en general y sobre c mo publicar en Anales AFA se encuentra disponible en la p gina de la revista. Este documento incluye las especificaciones de todos los estilos necesarios; puede por tanto ser utilizado como plantilla sobreescribiendo en el mismo. Por favor, NO INCLUYA numeraci n de p ginas.

\subsection{Tipos de letra y espaciado de l neas}

Todo manuscrito debe tener espaciado simple, increment ndolo autom ticamente para sub ndices y supra ndices. Debe usarse fuente Times New Roman y generarse en dos columnas. Se ahorran todos estos ajustes al usar esta plantilla en \LaTeX.
En la Tabla \ref{tab_01} se muestran los tama os de fuentes, formatos de p rrafo y separaci n entre secciones. Por ejemplo, el p rrafo normal requiere una sangr a de \unit[1]{cm} y una separaci n de 6 pts. entre p rrafos. Recuerde que 1 \emph{point} = 1 pt. $\cong$ \unit[0.25]{mm}. La lista de Estilos de la Tabla \ref{tab_01} tambi n est  incluida en el archivo asociado al presente documento.

\subsection{Tablas}

Las tablas deben colocarse lo m s cerca posible del punto donde se hace referencia por primera vez y deben citarse consecutivamente en orden num rico usando n meros ar bigos (1, 2, 3, etc.). 

Los t tulos de las tablas deben aparecer directamente arriba de las mismas.  Para generar una l nea horizontal gruesa en la tabla usar la opci n \color{Blue}\textsf{$\backslash$thickhline}\color{Black}. 

%EJEMPLO: TABLA QUE ABARCA 1 COLUMNA
\begin{table}[ht]
\caption{ Opciones del paquete \textsf{array} (entorno \textsf{tabular}).}
\label{tab_02}
    \centering
    \small           % modificaci n del tama o de la fuente
    \begin{tabular}{|m{5.5cm}|c|}
    \hline
    \centering
    \textbf{Descripci n} & \textbf{S mbolo}  \\ 
    \thickhline
    Alineaci n a la izquierda & l  \\ 
    \hline
    Alineaci n centrada & c  \\ 
    \hline
    Alineaci n a la derecha & r \\
    \hline
    Genera una caja del ancho indicado & p\{width\} \\
    \hline
    Inserta una l nea vertical entre columnas & $\mid$ \\
    \hline
    Genera una caja del ancho indicado. La entrada se centrar  verticalmente.  &  m\{width\}  \\
    \hline
    \end{tabular}
\end{table}

\normalsize     % fuente en tama o normal

El ancho de las tablas puede extenderse a las dos columnas. Para ello se recomienda usar la orden \color{Blue} \textsf{$\backslash$begin\{table*\} ... $\backslash$end\{table*\}} \color{Black}. En esta plantilla se presentan ejemplos de cada caso (ver Tablas \ref{tab_01} y \ref{tab_02}).


\subsection{Figuras}

Las figuras se incluyen en el documento en el sitio en que el autor lo considere m s conveniente. Numere cada figura consecutivamente con un n mero ar bigo de acuerdo con el orden en que se discute en el documento. Las leyendas de las figuras deben anotarse debajo de ellas (ver Fig. \ref{fig04}). 

\begin{figure}[ht]
\centering
   \includegraphics[width=0.48\textwidth]{fig2.jpg}
\caption{Modelo de figura en una columna.}
\label{fig04}
\end{figure} 

El ancho de las figuras tambi n puede extenderse a las dos columnas (ver Fig. \ref{fig05}).

Trate los gr ficos de modo que sean completamente autoexplicativos. Rotule cada eje (horizontal y vertical) incluyendo las correspondientes unidades, que deben estar separadas y encerradas entre corchetes (ejemplo: \textit{I (mA)}). Si es posible no use potencias de 10; en su lugar, use los prefijos apropiados del Sistema Internacional (ver Tabla \ref{tab_03}).
% EJEMPLO: FIGURA EN DOS COLUMNAS
\begin{figure*}[ht]
\centering
   \includegraphics[width=0.9\textwidth]{fig3.jpg}
\caption{Modelo de figura ocupando dos columnas.}
\label{fig05}
\end{figure*} 


\subsection{Formato de las ecuaciones}

Las ecuaciones deben centrarse en la columna y con su n mero de orden entre par ntesis, alineado por derecha. La numeraci n debe ser consecutiva usando n meros ar bigos. Se recomienda utilizar los entornos de ecuaciones disponibles en el paquete \textcolor{Blue}{\textsf{amsmath}},  recordando que
la orden \textcolor{Blue}{\textsf{$\backslash$begin$\{$equation$\}$}} numera autom ticamente las ecuaciones. Para escribir ecuaciones sin numeraci n utilice  \textcolor{Blue}{\textsf{$\backslash$begin$\{$equation*$\}$}}.
Trate el texto y las f rmulas matem ticas como una entidad. Punt e las expresiones matem ticas como oraciones o partes de oraciones. 
% ECUACIONES EN UNA L NEA 
Ejemplo: La soluci n es de la forma
\begin{equation}
r_{s} = A e^{\alpha^{+}s} + A e^{\alpha^{-}s}
\label{eq01}
\end{equation}
donde 
\begin{equation*}
e^{\alpha^{\pm}} = b \pm \sqrt{b^{2}+1}.
\label{eq01a}
\end{equation*}


Las expresiones matem ticas a menudo deben mostrarse en dos o m s l neas debido a limitaciones en el ancho de la columna definido por la revista. Las ecuaciones (\ref{eq02}) y (\ref{eq03}) son ejemplos de este caso. Para estos casos se recomienda utilizar los entornos  \\
\color{Blue}
\textsf{$\backslash$begin$\{$eqnarray$\}$} \\ \textsf{$\backslash$end$\{$eqnarray$\}$} \\
\color{black} 

% ECUACI N EN DOS L NEAS 
\begin{eqnarray}
\beta(\mathbf{k}) & \approx & \beta^{(0)}\,\delta(\mathbf{k}-\mathbf{k}_i) + \beta^{(1)}\,Z(\mathbf{k}) \nonumber\\
 & + & \int d^2k'\,\beta^{(2)}(\mathbf{k},\mathbf{k}')\,Z(\mathbf{k}-\mathbf{k}'),
\label{eq02}
\end{eqnarray}
para generar una ecuaci n en dos renglones    bien \\
\color{Blue}
\textsf{$\backslash$begin$\{$equation$\}$} \\ 
\textsf{$\backslash$begin$\{$split$\}$}  \\
\textsf{$\backslash$end$\{$split$\}$}  \\
\textsf{$\backslash$end$\{$equation$\}$}
\color{black}


% ECUACI N EN DOS L NEAS 
\begin{equation}
\begin{split}
& E_{x}=-\frac{1}{c}\frac{\partial A_{x}}{\partial t}-\frac{\partial \Phi}{\partial x} = - (\partial^{0}A^{1}-\partial^{1}A^{0}) \\
& B_{x}=\frac{\partial A_{z}}{\partial y}-\frac{\partial A_{y}}{\partial z}=-(\partial^{2}A^{3}-\partial^{3}A^{2}),
\end{split}
\label{eq03}
\end{equation}
para producir dos (o m s) ecuaciones etiquetadas con un s lo n mero.


Si la ecuaci n es muy larga, como en la ecuaci n (\ref{eq04}), y no es posible o apropiado mostrarla en dos o m s l neas, puede optar por presentarla en dos columnas. En este caso puede emplear un flotante extendido con el comando \textcolor{Blue}{\textsf{$\backslash$begin$\{$figure*$\}$}}
para posicionar la ecuaci n en el ancho total de la p gina.

% ECUACI N EN DOS COLUMNAS
\begin{figure*}
\begin{eqnarray}
[\mathbf{E}(\mathbf{\nabla . E}) - \mathbf{E}\times (\mathbf{\nabla \times E})]_{1}  & = & 
E_{1}\left( \frac{\partial E_{1}}{\partial x_{1}} +  \frac{\partial E_{2}}{\partial x_{2}} + \frac{\partial E_{3}}{\partial x_{3}} \right) 
- E_{2}\left( \frac{\partial E_{2}}{\partial x_{1}} -  \frac{\partial E_{1}}{\partial x_{2}} \right)
+ E_{3}\left( \frac{\partial E_{1}}{\partial x_{3}} -  \frac{\partial E_{3}}{\partial x_{1}} \right)
\nonumber\\
 & = & \frac{\partial}{\partial x_{1}}(E_{1}^{2}) + \frac{\partial}{\partial x_{2}}(E_{1}E_{2}) + \frac{\partial}{\partial x_{3}}(E_{1}E_{3}) - \frac{1}{2}\frac{\partial}{\partial x_{1}}(E_{1}^{2}+E_{2}^{2}+E_{3}^{2})
\label{eq04}
\end{eqnarray}
\end{figure*}

\subsection{Agrupaci n de s mbolos: jerarqu a}

Con el prop sito de agrupar, la secuencia de corchetes preferida para los art culos publicados en Anales AFA es \large{\{[()]\}}. \normalsize Ejemplo:
\begin{equation}
D=3\left[L^{2}+\frac{R}{(L+1.5)^{2}}\right]+(R-3L)R. 
\label{eq05}
\end{equation}

\subsection{Unidades}

Utilice las reglas y convenciones aceptadas internacionalmente, como el Sistema Internacional de unidades (SI). Para acceder a una lista de s mbolos est ndar para unidades de medida  puede  consultar en \urlstyle{same}\url{https://doi.org/10.6028/NIST.SP.330-2019}.

Los m ltiplos decimales o subm ltiplos de unidades se indican mediante el uso de prefijos como se muestra en la Tabla \ref{tab_03}).

\begin{table}[ht]
\caption{Prefijos del Sistema Internacional. }
\label{tab_03}
    \centering
    \small               % modificaci n del tama o de la fuente
    \begin{tabular}{|c|c|c|c|c|c|}
    \hline
    \centering
    \textbf{Prefijo} & \textbf{S mbolo} & \textbf{Factor} & \textbf{Prefijo} & \textbf{S mbolo} & \textbf{Factor}   \\ 
    \thickhline
    deci & d & $10^{-1}$ & deca & da &  $10^{1}$ \\ 
    \hline
    centi & c & $10^{-2}$ & hecto & h &  $10^{2}$  \\ 
    \hline
    mili & m & $10^{-3}$ & kilo & k &  $10^{3}$  \\
    \hline
    micro & $\mu$ & $10^{-6}$ & mega & M &  $10^{6}$  \\
    \hline
    nano & n & $10^{-9}$ & giga & G &  $10^{9}$  \\
    \hline
    pico & p & $10^{-12}$ & tera & T &  $10^{12}$   \\
    \hline
    femto & f & $10^{-15}$ & peta & P &  $10^{15}$  \\
    \hline
    atto & a & $10^{-18}$ & exa & E &  $10^{18}$   \\
    \hline
    \end{tabular}
\end{table}

\normalsize

Para contemplar el espacio entre el valor de la magnitud y la unidad se recomienda usar el comando 
\textcolor{Blue}{\textsf{$\backslash$unit[VALOR]\{UNIDAD\}}}. Por ejemplo, \textcolor{Blue}{\textsf{$\backslash$unit[1,2]\{cm\}}}  da como resultado \unit[1,2]{cm} (no escribir 1,2cm). 

Tenga en cuenta que algunas unidades no est n separadas del valor de la magnitud. Ejemplos: 1$\%$, 1  y 1 C.


\subsection{Formato de las referencias} 
Para las referencias se sigue el formato de las revistas de la APS. 
Citar las referencias utilizando n meros ar bigos consecutivos entre corchetes por orden de aparici n en el manuscrito. 

En caso de generar sus archivos en \LaTeX{} se recomienda a los autores que utilicen BibTEX para la generaci n de las entradas bibliogr ficas. Si se utiliza BibTEX, durante el proceso de producci n final se requerir  el archivo .bib  junto con el archivo .tex principal del manuscrito.

Al usar esta plantilla, el listado de referencias se genera autom ticamente por orden de aparici n en el texto. 
El comando \color{Blue} \textsf{$\backslash$cite} \color{Black} producira un n mero de referencia de la forma [\color{blue}N MERO\color{Black}]. 

En el listado de referencias, incluir los t tulos de los art culos, libros, monograf as, tesis, etc. La abreviaci n de revistas puede consultarse en \urlstyle{same}\url{https://images.webofknowledge.com/WOK46P9/help/WOS/A_abrvjt.html}. 


Datos bibliogr ficos necesarios para la elaboraci n de la referencia:

\begin{itemize}[itemsep=1mm, parsep=0pt]
    \item Libros (@book): Autores. T tulo del libro. Editorial. Lugar de edici n. Nro. de edici n. A o de publicaci n. Ver ejemplo en \cite{libro1}.
    \item Art culos de revistas peri dicas (@article): Autores. T tulo del trabajo. Nombre de la revista. Volumen. N mero. P ginas. A o de publicaci n. Ver ejemplos en \cite{paper1} y \cite{paper2}.
    \item Art culos de comunicaciones a congresos (@inproceedings): Autores. T tulo del trabajo. T tulo del libro. P ginas donde aparece publicada la contribuci n. Lugar de edici n. Fecha de publicaci n. Editorial. Ver ejemplo en \cite{proceeding1}.
    \item Tesis doctorales (@phdthesis): Autor. T tulo de la tesis. Facultad, Universidad. A o. Ver ejemplo en \cite{tesis1}.
    \item Tesis de maestr a/licenciatura (@mastersthesis): Autor. T tulo del trabajo. Facultad, Universidad. A o. Ver ejemplo en \cite{maestria1}.
    \item Material variado (@misc): T tulo del documento. Direcci n electr nica (en caso que corresponda). Fecha de consulta. Opcional: autor. Ver ejemplo en \cite{motor}.
    \item Informe publicado por alguna instituci n (@techreport): Autores. T tulo del informe. Instituci n. A o. Ver ejemplo en \cite{informe1}.
\end{itemize}


%\subsection{Material ilustrativo}
% Las URL deben especificarse as : \urlstyle{same}\url{www.patagonia.com.ar}) 

%-------------------------------------------------------------
% BIBLIOGRAF A (edici n: AUTORES. Ir a referencias.bib)
\section*{REFERENCIAS} 
 \printbibliography[heading=none] 

%-------------------------------------------------------------
% Fin del documento
\end{document}